\documentclass[a4paper,oneside,12pt]{report}
% Package yang diperlukan
\usepackage[table]{xcolor} % Untuk menggunakan warna
\usepackage{graphicx}      % Untuk memasukkan gambar
\usepackage{titling}       % Untuk kustomisasi halaman judul
\usepackage{geometry}      % Untuk mengatur margin halaman
\usepackage{listings}      % Untuk menampilkan blok kode
\usepackage{float}         % Untuk kontrol penempatan float
\usepackage{hyperref}      % Untuk membuat hyperlink
\usepackage{indentfirst}   % Untuk memberi indentasi pada paragraf pertama
\usepackage{url}           % Untuk format URL
\usepackage{fancyhdr}      % Untuk kustomisasi header dan footer
\usepackage{tocloft}       % Untuk kustomisasi Daftar Isi
\usepackage{bookmark}      % Untuk manajemen bookmark PDF
\usepackage{amsmath}       % Untuk notasi matematika
\usepackage[most]{tcolorbox} % Untuk box output terminal

% Paket untuk tabel
\usepackage{array}
\usepackage{booktabs}
\usepackage{tabularx}
\usepackage{multirow}
\usepackage{colortbl}

% Macro dan warna untuk tabel
\newcommand{\thead}[1]{\textcolor{white}{\textbf{#1}}}
\definecolor{tableheader}{RGB}{30, 58, 138}
\definecolor{tablezebra}{RGB}{248, 250, 252}
\definecolor{tableborder}{RGB}{203, 213, 225}

%======================================================================
% PENGATURAN GEOMETRI HALAMAN
%======================================================================
\geometry{left=3cm, right=2.5cm, top=2cm, bottom=2.5cm}
\sloppy

%======================================================================
% PENGATURAN WARNA UNTUK BLOK KODE
%======================================================================
\definecolor{codegreen}{rgb}{0,0.6,0}
\definecolor{codegray}{rgb}{0.5,0.5,0.5}
\definecolor{codepurple}{rgb}{0.58,0,0.82}
\definecolor{backcolour}{rgb}{0.95,0.95,0.92}
\definecolor{darkblue}{rgb}{0,0,0.6}
\definecolor{outputbg}{RGB}{241,245,249}
\definecolor{outputborder}{RGB}{203,213,225}

%======================================================================
% PENGATURAN BLOK KODE (LISTINGS)
%======================================================================
\lstdefinestyle{commonstyle}{
    backgroundcolor=\color{backcolour},
    breakatwhitespace=false,
    breaklines=true,
    captionpos=b,
    keepspaces=true,
    numbers=left,
    numbersep=5pt,
    showspaces=false,
    showstringspaces=false,
    showtabs=false,
    tabsize=2,
    frame=single,
    basicstyle=\ttfamily\footnotesize,
    numberstyle=\tiny\color{codegray},
    columns=flexible,
    postbreak=\mbox{\textcolor{red}{$\hookrightarrow$}\space},
    breakindent=0pt,
    breakautoindent=false
}

\lstdefinestyle{pythonstyle}{
    style=commonstyle,
    language=Python,
    commentstyle=\color{codegreen},
    keywordstyle=\color{darkblue},
    stringstyle=\color{codepurple},
    emph={import,from,class,def,for,while,if,else,elif,try,except,finally,as,lambda,return},
    emphstyle=\color{darkblue}
}

% Macro untuk gambar tangkapan layar Colab
\newcommand{\screenshotimage}[4][0.88\textwidth]{%
    \begin{figure}[H]
        \centering
        \includegraphics[width=#1]{#2}
        \caption{#3}
        \label{fig:#4}
    \end{figure}
}

%======================================================================
% PENGATURAN LAIN-LAIN
%======================================================================
\renewcommand{\figurename}{Gambar}
\renewcommand{\thefigure}{\arabic{figure}}
\renewcommand{\tablename}{Tabel}

\hypersetup{
    colorlinks=true,
    linkcolor=black,
    filecolor=black,
    urlcolor=black,
    citecolor=black,
    pdftitle={Laporan Praktikum Penambangan Data Pertemuan VIII},
    pdfpagemode=FullScreen,
    breaklinks=true
}

\pagestyle{fancy}
\fancyhf{}
\fancyfoot[C]{\thepage}
\renewcommand{\headrulewidth}{0pt}

\fancypagestyle{plain}{
    \fancyhf{}
    \fancyfoot[C]{\thepage}
    \renewcommand{\headrulewidth}{0pt}
}

%======================================================================
% PENGATURAN STRUKTUR BAB DAN DAFTAR ISI
%======================================================================
\renewcommand{\contentsname}{Daftar Isi}
\renewcommand{\bibname}{Daftar Pustaka}

\renewcommand{\thechapter}{}
\renewcommand{\chaptername}{}

\renewcommand{\thesection}{\arabic{chapter}.\arabic{section}}
\renewcommand{\thesubsection}{\thesection.\arabic{subsection}}

\setcounter{secnumdepth}{2}
\setcounter{tocdepth}{2}

\renewcommand{\cftchapdotsep}{\cftdotsep}
\renewcommand{\cftchapleader}{\cftdotfill{\cftchapdotsep}}
\setlength{\cftchapindent}{0em}
\setlength{\cftsecindent}{2em}
\setlength{\cftsubsecindent}{4em}
\setlength{\cftchapnumwidth}{0em}
\setlength{\cftsecnumwidth}{2.5em}
\setlength{\cftsubsecnumwidth}{3.2em}
\renewcommand{\cftchappresnum}{}
\renewcommand{\cftchapaftersnum}{}

%======================================================================
% DOKUMEN UTAMA
%======================================================================
\begin{document}

%======================================================================
% HALAMAN JUDUL (COVER)
%======================================================================
\newgeometry{left=3cm, right=2.5cm, top=2.5cm, bottom=3cm}
\begin{titlingpage}
\begin{center}
\vspace*{0.5cm}
{\large \textbf{\MakeUppercase{Laporan Praktikum}}} \\[0.4cm]
{\large \textbf{\MakeUppercase{Praktikum Penambangan Data}}} \\[0.4cm]
{\large \textbf{\MakeUppercase{Pertemuan VIII}}} \\[0.4cm]
{\large \textbf{\MakeUppercase{Neural Network (Multilayer Perceptron)}}} \\[1.5cm]

\includegraphics[height=7cm]{lambang ugm.png} \\[1cm]

{\normalsize \textbf{Disusun oleh:}} \\[0.4cm]
\begin{tabular}{@{}l l @{}}
  \textbf{Nama}           & : Januarsyah Akbar \\
  \textbf{NIM}            & : 24/535846/SV/24314 \\
  \textbf{Kelas}          & : B2 \\
  \textbf{Dosen Pengampu} & : Dr. Ganjar Alfian, S.T., M.Eng. \\
                          & \phantom{:} Dr. Imam Fahrurrozi, S.T., M.Cs.
\end{tabular} \\[1.2cm]

{\normalsize \textbf{PROGRAM STUDI D-IV}} \\[0.2cm]
{\normalsize \textbf{TEKNOLOGI REKAYASA PERANGKAT LUNAK}} \\[0.2cm]
{\normalsize \textbf{DEPARTEMEN TEKNIK ELEKTRO DAN INFORMATIKA}} \\[0.2cm]
{\normalsize \textbf{SEKOLAH VOKASI}} \\[0.2cm]
{\normalsize \textbf{UNIVERSITAS GADJAH MADA}} \\[0.2cm]
{\normalsize \textbf{YOGYAKARTA}} \\[0.2cm]
{\normalsize \textbf{2026}} \\[1.2cm]
\end{center}
\end{titlingpage}
\restoregeometry

\thispagestyle{empty}
\clearpage
\stepcounter{page}
\thispagestyle{empty}

% --- Daftar Isi ---
\addcontentsline{toc}{chapter}{Daftar Isi}
\tableofcontents
\clearpage

\pagenumbering{arabic}
\setcounter{page}{3}

%======================================================================
% BAB I: TUJUAN PRAKTIKUM
%======================================================================
\chapter[Tujuan Praktikum]{Tujuan Praktikum}

\begin{enumerate}
    \item Mahasiswa memahami konsep Multilayer Perceptron (MLP) dan mampu menerapkannya pada teknik klasifikasi serta regresi.
    \item Mahasiswa mampu menganalisis hasil klasifikasi menggunakan MLP dengan memanfaatkan metrik evaluasi dan pengujian performa model.
    \item Mahasiswa mampu menganalisis hasil regresi menggunakan MLP dengan membandingkan nilai hasil prediksi terhadap data asli.
\end{enumerate}

\section*{Alat dan Bahan}
\begin{itemize}
    \item \textbf{Perangkat Keras}:
    \begin{enumerate}
        \item Laptop dengan kapasitas RAM minimal 4 GB.
        \item Koneksi internet stabil untuk akses repositori dan komputasi awan.
    \end{enumerate}
    \item \textbf{Perangkat Lunak}:
    \begin{enumerate}
        \item Peramban web modern (Google Chrome, Mozilla Firefox, Microsoft Edge, atau Safari).
        \item Lingkungan Google Colab dan Jupyter Notebook berbasis Python 3.11.
        \item Pustaka penambangan data dan pembelajaran mesin (\texttt{pandas}, \texttt{numpy}, \texttt{matplotlib}, \texttt{seaborn}, \texttt{scikit-learn}, dan \texttt{imbalanced-learn}).
    \end{enumerate}
\end{itemize}

\clearpage

%======================================================================
% BAB II: DASAR TEORI
%======================================================================
\chapter[Dasar Teori]{Dasar Teori}

\section{Neural Network}
\textit{Neural Network} atau jaringan syaraf tiruan adalah arsitektur komputasi yang terinspirasi oleh sistem jaringan saraf biologis otak manusia \cite{ref_haykin}. Jaringan ini tersusun atas interkoneksi unit pemrosesan dasar yang disebut neuron buatan (\textit{perceptron}). Setiap neuron menerima sinyal masukan dari lapisan sebelumnya, mengalikannya dengan bobot sinaptik tertentu, menjumlahkannya dengan nilai bias, lalu meneruskan hasil akumulasi melalui suatu fungsi aktivasi menuju lapisan berikutnya \cite{ref_goodfellow}.

Struktur jaringan syaraf tiruan dirancang untuk memodelkan fungsi pemetaan matematis yang kompleks dan non-linear \cite{ref_hukubun}. Melalui proses pelatihan terbimbing (\textit{supervised learning}), jaringan memperbarui bobot secara bertahap berdasarkan umpan balik kesalahan keluaran. Keunggulan utama neural network terletak pada fleksibilitas representasinya, mampu menangani volume data masukan berskala masif, mengenali korelasi implisit antardimensi, serta melakukan generalisasi pola data baru secara adaptif \cite{ref_nielsen}.

\section{Multilayer Perceptron (MLP)}
\textit{Multilayer Perceptron} (MLP) merupakan varian arsitektur jaringan syaraf tiruan umpan maju (\textit{feedforward artificial neural network}) yang terdiri dari minimal tiga lapisan neuron: satu lapisan masukan (\textit{input layer}), satu atau lebih lapisan tersembunyi (\textit{hidden layer}), dan satu lapisan keluaran (\textit{output layer}) \cite{ref_pardede}. Keberadaan lapisan tersembunyi memungkinkan jaringan memecahkan persoalan yang tidak dapat dipisahkan secara linier (\textit{non-linearly separable problems}), seperti masalah logika XOR yang gagal diselesaikan oleh perseptron tunggal \cite{ref_haykin}.

Secara matematis, keluaran neuron ke-$j$ pada lapisan tersembunyi dirumuskan sebagai berikut:
\begin{equation}
    z_j = f \left( \sum_{i=1}^{n} w_{ji} x_i + b_j \right)
\end{equation}
di mana $x_i$ menyatakan nilai masukan ke-$i$, $w_{ji}$ adalah bobot sinaptik antara masukan $i$ dan neuron $j$, $b_j$ merupakan bias neuron $j$, dan $f(\cdot)$ adalah fungsi aktivasi non-linear. Keunggulan MLP meliputi ketahanan terhadap derau data, kapabilitas estimasi probabilitas posterior kelas, serta fleksibilitas penggunaan untuk permasalahan klasifikasi maupun estimasi nilai kontinu pada regresi \cite{ref_pardede}.

\section{Backpropagation Learning}
\textit{Backpropagation} adalah algoritma pembelajaran standar yang digunakan untuk mengoptimalkan kumpulan bobot dan bias pada Multilayer Perceptron \cite{ref_goodfellow}. Mekanisme kerja algoritma ini terbagi dalam dua fase utama: perambatan maju (\textit{forward pass}) dan perambatan mundur (\textit{backward pass}). Pada fase maju, sinyal masukan diproses lapis demi lapis hingga menghasilkan nilai prediksi akhir pada lapisan keluaran, yang kemudian dibandingkan dengan label target aktual untuk menghitung nilai fungsi kerugian (\textit{loss function}) $E$ \cite{ref_haykin}.

Pada fase perambatan mundur, sinyal kesalahan disebarkan kembali dari lapisan keluaran menuju lapisan masukan. Gradien parsial dari fungsi kerugian terhadap masing-masing bobot dihitung menggunakan aturan rantai kalkulus (\textit{chain rule}):
\begin{equation}
    \frac{\partial E}{\partial w_{ji}} = \frac{\partial E}{\partial z_j} \cdot \frac{\partial z_j}{\partial \text{net}_j} \cdot \frac{\partial \text{net}_j}{\partial w_{ji}}
\end{equation}
Bobot kemudian diperbarui menggunakan metode penurunan gradien (\textit{gradient descent}) atau optimizer adaptif seperti Adam:
\begin{equation}
    w_{ji}^{(\text{baru})} = w_{ji}^{(\text{lama})} - \eta \frac{\partial E}{\partial w_{ji}}
\end{equation}
dengan $\eta$ melambangkan laju pembelajaran (\textit{learning rate}). Proses iteratif ini berlangsung terus hingga model mencapai titik konvergensi minimum lokal atau global \cite{ref_goodfellow}.

\section{Activation Function}
Fungsi aktivasi memegang peranan krusial dalam neural network karena menyuntikkan sifat non-linearitas ke dalam sistem komputasi \cite{ref_nielsen}. Tanpa fungsi aktivasi non-linear, operasi beruntun pada jaringan multi-lapis hanya akan menghasilkan kombinasi linier sederhana yang ekuivalen dengan model regresi linier satu lapis \cite{ref_haykin}. Ragam fungsi aktivasi yang umum diimplementasikan meliputi:
\begin{enumerate}
    \item \textbf{Sigmoid (Logistic)}: Memetakan nilai riil ke rentang $(0, 1)$, sering digunakan pada probabilitas biner.
    \begin{equation}
        \sigma(x) = \frac{1}{1 + e^{-x}}
    \end{equation}
    \item \textbf{Hyperbolic Tangent (Tanh)}: Memetakan masukan ke rentang $(-1, 1)$, memiliki sifat nol terpusat (\textit{zero-centered}).
    \begin{equation}
        \tanh(x) = \frac{e^x - e^{-x}}{e^x + e^{-x}}
    \end{equation}
    \item \textbf{Rectified Linear Unit (ReLU)}: Mempertahankan nilai positif dan menolkan nilai negatif, efisien secara komputasi serta meminimalkan fenomena peluruhan gradien (\textit{vanishing gradient}).
    \begin{equation}
        \text{ReLU}(x) = \max(0, x)
    \end{equation}
    \item \textbf{Softmax}: Digunakan pada lapisan keluaran klasifikasi multi-kelas untuk menghasilkan distribusi probabilitas normal dengan total penjumlahan bernilai satu.
    \begin{equation}
        \text{Softmax}(x_i) = \frac{e^{x_i}}{\sum_{j=1}^{K} e^{x_j}}
    \end{equation}
\end{enumerate}

\section{Klasifikasi}
Klasifikasi adalah paradigma pembelajaran terarah (\textit{supervised learning}) yang bertujuan memprediksi kategori diskret suatu data observasi berdasarkan himpunan fitur terukur \cite{ref_roihan}. Model klasifikasi dilatih menggunakan pasangan fitur dan label kelas aktual. Ragam algoritma klasifikasi mencakup Support Vector Machines, Naive Bayes, Decision Tree, Random Forest, serta Multilayer Perceptron (\texttt{MLPClassifier}) \cite{ref_roihan}.

Evaluasi model klasifikasi diukur menggunakan \textit{Confusion Matrix} yang memetakan empat kondisi prediksi: True Positive (TP), True Negative (TN), False Positive (FP), dan False Negative (FN). Metrik turunan utama meliputi:
\begin{equation}
    \text{Accuracy} = \frac{\text{TP} + \text{TN}}{\text{TP} + \text{TN} + \text{FP} + \text{FN}}
\end{equation}
\begin{equation}
    \text{Precision} = \frac{\text{TP}}{\text{TP} + \text{FP}}
\end{equation}
\begin{equation}
    \text{Recall} = \frac{\text{TP}}{\text{TP} + \text{FN}}
\end{equation}
\begin{equation}
    F_1\text{-Score} = 2 \times \frac{\text{Precision} \times \text{Recall}}{\text{Precision} + \text{Recall}}
\end{equation}

\section{Regresi}
Regresi adalah teknik pemodelan prediktif yang berfokus pada estimasi nilai kuantitatif kontinu (\textit{continuous numerical values}) dari variabel dependen berdasarkan satu atau lebih variabel independen \cite{ref_mardiani}. Berbeda dari regresi linier parametrik klasik yang berasumsi adanya hubungan proporsional linier kaku, Multilayer Perceptron Regressor (\texttt{MLPRegressor}) mampu menangkap pola non-linear rumit tanpa memerlukan perumusan eksplisit bentuk fungsi matematis di awal \cite{ref_haykin}.

Pengukuran akurasi dan kinerja model regresi dievaluasi menggunakan metrik berikut:
\begin{enumerate}
    \item \textbf{Koefisien Determinasi ($R^2$)}: Mengukur proporsi variansi data target yang berhasil dijelaskan oleh fitur masukan model.
    \begin{equation}
        R^2 = 1 - \frac{\sum_{i=1}^{n} (y_i - \hat{y}_i)^2}{\sum_{i=1}^{n} (y_i - \bar{y})^2}
    \end{equation}
    \item \textbf{Mean Squared Error (MSE)}: Rata-rata kuadrat selisih error antara nilai aktual $y_i$ dan nilai estimasi $\hat{y}_i$.
    \begin{equation}
        \text{MSE} = \frac{1}{n} \sum_{i=1}^{n} (y_i - \hat{y}_i)^2
    \end{equation}
    \item \textbf{Root Mean Squared Error (RMSE)}: Akar kuadrat dari nilai MSE, memiliki satuan besaran yang sama dengan variabel target asli sehingga mudah diinterpretasikan.
    \begin{equation}
        \text{RMSE} = \sqrt{\frac{1}{n} \sum_{i=1}^{n} (y_i - \hat{y}_i)^2}
    \end{equation}
\end{enumerate}

\clearpage

%======================================================================
% BAB III: HASIL DAN PEMBAHASAN
%======================================================================
\chapter[Hasil dan Pembahasan]{Hasil dan Pembahasan}

\section{Langkah Praktikum Percobaan}

\subsection{Persiapan Lingkungan Komputasi}
Eksperimen dilaksanakan pada lingkungan Google Colab dan Jupyter Notebook lokal dengan runtime Python 3.11. Lingkungan ini memfasilitasi komputasi matriks berkecepatan tinggi serta menyediakan integrasi pustaka pengolahan data dan pembelajaran mesin tanpa beban konfigurasi perangkat keras fisik yang kompleks.

\subsection{Instalasi dan Impor Library}
Tahap awal percobaan diawali dengan mengimpor modul manipulasi data, pra-pemrosesan fitur, fungsi pembagian data latih dan uji, modul arsitektur neural network, serta ragam metrik evaluasi kinerja.

\begin{lstlisting}[style=pythonstyle, caption={Impor Pustaka Pembelajaran Mesin dan MLP}]
from sklearn.model_selection import train_test_split
from sklearn.preprocessing import LabelEncoder
from sklearn.metrics import accuracy_score
from sklearn.metrics import precision_score
from sklearn.metrics import recall_score
from sklearn.metrics import confusion_matrix
from sklearn.metrics import ConfusionMatrixDisplay
from imblearn.metrics import sensitivity_specificity_support
from sklearn.preprocessing import StandardScaler
from sklearn.preprocessing import OrdinalEncoder
from sklearn.neural_network import MLPClassifier

import os
import pandas as pd
import numpy as np
import matplotlib.pyplot as plt
\end{lstlisting}

\screenshotimage[0.92\textwidth]{gambar/colab-p8-imports.png}{Impor Pustaka Preprocessing, Model MLP, dan Metrik Evaluasi (Gambar 8.4.1 Modul)}{colab-imports}

\subsection{Import dan Load Dataset Information}
Dataset yang digunakan dalam langkah percobaan terdiri atas dua jenis domain:
\begin{enumerate}
    \item \textbf{Klasifikasi}: Dataset Bank Customer Churn (\texttt{BankChurners.csv}), berisikan 10.127 data transaksi dan profil nasabah perbankan untuk memprediksi risiko penghentian layanan kartu kredit.
    \item \textbf{Regresi}: Dataset King County House Sales (\texttt{kc\_house\_data.csv}), berisikan 21.613 catatan penjualan rumah residensial untuk mengestimasi nilai harga properti.
\end{enumerate}

\subsection{Percobaan Klasifikasi dengan MLPClassifier}

\subsubsection{Membaca Dataset Bank Customer Churn}
Dataset dibaca menggunakan fungsi \texttt{pd.read\_csv()}. Lima baris teratas diperiksa melalui perintah \texttt{df.head()} untuk mengamati atribut awal.

\begin{lstlisting}[style=pythonstyle, caption={Pembacaan Dataset Bank Customer Churn}]
# Membaca dataset
df = pd.read_csv('BankChurners.csv')

# Menampilkan 5 baris data teratas
df.head()
\end{lstlisting}

\screenshotimage[0.95\textwidth]{gambar/colab-p8-churn-head.png}{Tampilan Lima Baris Pertama Dataset Bank Customer Churn (Gambar 8.4.2 Modul)}{colab-churn-head}

\subsubsection{Analisis Distribusi Variabel dengan Histogram}
Distribusi frekuensi seluruh fitur numerik dipetakan menggunakan fungsi \texttt{df.hist(figsize=(20, 15))}.

\begin{lstlisting}[style=pythonstyle, caption={Visualisasi Histogram Fitur Bank Churn}]
df.hist(figsize=(20, 15))
plt.show()
\end{lstlisting}

\screenshotimage[0.92\textwidth]{gambar/colab-p8-churn-hist.png}{Distribusi Histogram Fitur Bank Customer Churn (Gambar 8.4.3 Modul)}{colab-churn-hist}

\subsubsection{Pemeriksaan Proporsi Target Class}
Proporsi target klasifikasi \texttt{Attrition\_Flag} dihitung untuk meninjau keseimbangan kelas.

\begin{lstlisting}[style=pythonstyle, caption={Pemeriksaan Proporsi Target Attrition\_Flag}]
# Cek persentase target class
data = df['Attrition_Flag'].value_counts()
data.plot(kind='pie', autopct='%1.2f%%')
plt.show()
\end{lstlisting}

\screenshotimage[0.65\textwidth]{gambar/colab-p8-churn-pie.png}{Proporsi Distribusi Kelas Target Attrition\_Flag (Gambar 8.4.4 Modul)}{colab-churn-pie}

Berdasarkan grafik Gambar \ref{fig:colab-churn-pie}, kelas mayoritas nasabah aktif (\textit{Existing Customer}) mendominasi sebesar 83,93\% (8.500 nasabah), sedangkan kelas minoritas nasabah churn (\textit{Attrited Customer}) hanya mencakup 16,07\% (1.627 nasabah).

\subsubsection{Pemeriksaan Missing Value}
Pemeriksaan nilai kosong dilakukan dengan memanggil \texttt{df.isnull().sum()}.

\begin{lstlisting}[style=pythonstyle, caption={Pemeriksaan Nilai Kosong pada Bank Churn}]
# Pengecekan missing value
df.isnull().sum()
\end{lstlisting}

\screenshotimage[0.75\textwidth]{gambar/colab-p8-churn-null.png}{Hasil Pengecekan Missing Value pada Bank Churn (Gambar 8.4.5 Modul)}{colab-churn-null}

Hasil pemeriksaan menunjukkan seluruh atribut memiliki nilai 0, yang menandakan tidak terdapat nilai kosong (\textit{missing value}) pada dataset ini.

\subsubsection{Identifikasi Atribut Kategorikal dan Numerik}
Atribut non-numerik disaring menggunakan \texttt{select\_dtypes}. Ditemukan lima kolom bertipe kategorikal: \texttt{Gender}, \texttt{Education\_Level}, \texttt{Marital\_Status}, \texttt{Income\_Category}, dan \texttt{Card\_Category}.

\screenshotimage[0.85\textwidth]{gambar/colab-p8-churn-cats.png}{Identifikasi Atribut Kategorikal pada Dataframe Fitur (Gambar 8.4.6 Modul)}{colab-churn-cats}

\subsubsection{Pembagian Data Train-Test dan Pelatihan Model MLPClassifier}
Variabel target \texttt{Attrition\_Flag} dikonversi ke nilai biner (0 untuk Attrited Customer, 1 untuk Existing Customer) menggunakan \texttt{LabelEncoder}. Seluruh kolom kategorikal pada matriks fitur $X$ juga diubah ke representasi integer. Data dibagi menjadi 70\% porsi pelatihan dan 30\% porsi pengujian (\texttt{test\_size=0.3}, \texttt{random\_state=42}). Selanjutnya dilakukan standarisasi fitur menggunakan \texttt{StandardScaler} sebelum dialirkan ke \texttt{MLPClassifier}.

\begin{lstlisting}[style=pythonstyle, caption={Preprocessing dan Pelatihan MLPClassifier}]
# Menyiapkan X dan y
df_X = df.drop(['CLIENTNUM', 'Attrition_Flag'], axis=1)
df_y = df[['Attrition_Flag']]

# Label encoding target y
le = LabelEncoder()
df_y = le.fit_transform(df_y['Attrition_Flag'])

# Label encoding atribut kategorikal
cats = df_X.select_dtypes(include=['object', 'bool']).columns
for i in list(cats.values):
    le_col = LabelEncoder()
    df_X[i] = le_col.fit_transform(df_X[i])

X = df_X.astype(float).values
y = df_y.astype(float)

# Pembagian data latih dan uji
X_train, X_test, y_train, y_test = train_test_split(X, y, test_size=0.3, random_state=42)

# Standarisasi fitur
scaler = StandardScaler().fit(X_train)
X_train = scaler.transform(X_train)
X_test = scaler.transform(X_test)

# Pelatihan model MLP
model = MLPClassifier(max_iter=500, random_state=42)
model.fit(X_train, y_train)
\end{lstlisting}

\screenshotimage[0.90\textwidth]{gambar/colab-p8-churn-split-train.png}{Alur Pembagian Data, Standarisasi, dan Pelatihan Model (Gambar 8.4.7 Modul)}{colab-churn-split}

\subsubsection{Inspeksi Bobot dan Data Terproses}
Koefisien bobot model diperiksa melalui atribut \texttt{model.coefs\_}. Lapisan tersembunyi berukuran 100 neuron menghasilkan matriik bobot pertama berdimensi $(19, 100)$ yang menghubungkan 19 fitur masukan ke 100 unit tersembunyi, serta matriks kedua berdimensi $(100, 1)$ menuju neuron keluaran.

\screenshotimage[0.90\textwidth]{gambar/colab-p8-churn-coefs.png}{Struktur Koefisien Bobot Saraf Tiruan pada Model (Gambar 8.4.8 Modul)}{colab-churn-coefs}

\screenshotimage[0.95\textwidth]{gambar/colab-p8-churn-dfx.png}{Tampilan Dataframe Fitur Setelah Encoding (Gambar 8.4.9 Modul)}{colab-churn-dfx}

\screenshotimage[0.88\textwidth]{gambar/colab-p8-churn-xtrain0.png}{Representasi Vektor Numerik Sampel Pertama X\_train[0] (Gambar 8.4.10 Modul)}{colab-churn-xtrain0}

\subsubsection{Evaluasi Kinerja Model Klasifikasi}
Kinerja model diuji terhadap 3.039 data uji independen. Metrik evaluasi mencakup Akurasi, Presisi, Recall, dan visualisasi \textit{Confusion Matrix}.

\begin{lstlisting}[style=pythonstyle, caption={Evaluasi Metrik dan Confusion Matrix}]
y_pred = model.predict(X_test)
print("Accuracy :", accuracy_score(y_test, y_pred))
print("Precision :", precision_score(y_test, y_pred))
print("Recall :", recall_score(y_test, y_pred))
print("Confusion matrix :\n", confusion_matrix(y_test, y_pred))

disp = ConfusionMatrixDisplay(confusion_matrix=confusion_matrix(y_test, y_pred), display_labels=model.classes_)
disp.plot()
plt.show()
\end{lstlisting}

\screenshotimage[0.75\textwidth]{gambar/colab-p8-churn-eval-metrics.png}{Ringkasan Output Numerik Metrik Evaluasi Klasifikasi}{colab-churn-eval}

\screenshotimage[0.65\textwidth]{gambar/colab-p8-churn-cm.png}{Visualisasi Confusion Matrix MLPClassifier (Gambar 8.4.11 Modul)}{colab-churn-cm}

Dari hasil evaluasi, model \texttt{MLPClassifier} mencatatkan akurasi sebesar \textbf{93,81\%}, presisi sebesar \textbf{96,05\%}, dan recall sebesar \textbf{96,58\%}. Dari matriks kontingensi, model berhasil mengklasifikasikan 2.456 nasabah aktif dengan tepat dan mendeteksi 395 nasabah yang berpotensi churn.

\subsubsection{Prediksi Data Inferensi Baru}
Satu sampel nasabah baru dengan profil usia 58 tahun, perempuan, berpendidikan Graduate, status lajang, pendapatan di bawah \$40K, kartu tipe Blue, saldo revolving \$800, dan utilisasi kredit 16,7\% diumpankan ke dalam model.

\screenshotimage[0.90\textwidth]{gambar/colab-p8-churn-inference-code.png}{Penyiapan Dataframe Inferensi Nasabah Baru (Gambar 8.4.12 Modul)}{colab-churn-inf-code}

\screenshotimage[0.85\textwidth]{gambar/colab-p8-churn-inference-out.png}{Hasil Prediksi Inferensi Model Terhadap Nasabah Baru (Gambar 8.4.13 Modul)}{colab-churn-inf-out}

Model memprediksi keluaran kelas 1, yang diinterpretasikan sebagai \textbf{Existing Customer} dengan tingkat probabilitas keyakinan mencapai 96,52\%.

\subsection{Percobaan Regresi dengan MLPRegressor}

\subsubsection{Membaca Dataset KC House Data}
Dataset harga properti King County dibaca menggunakan \texttt{pd.read\_csv()}.

\screenshotimage[0.95\textwidth]{gambar/colab-p8-kc-head.png}{Lima Baris Teratas King County House Sales Data (Gambar 8.4.14 Modul)}{colab-kc-head}

\subsubsection{Pemeriksaan Missing Value dan Distribusi Numerik}
Pemeriksaan \texttt{df.isnull().sum()} membuktikan bahwa seluruh kolom bersih dari nilai kosong. Distribusi nilai fitur divisualisasikan dalam bentuk histogram.

\screenshotimage[0.75\textwidth]{gambar/colab-p8-kc-null.png}{Pengecekan Missing Value KC House Data (Gambar 8.4.15 Modul)}{colab-kc-null}

\screenshotimage[0.92\textwidth]{gambar/colab-p8-kc-hist.png}{Distribusi Histogram Fitur Numerik KC House Data (Gambar 8.4.16 Modul)}{colab-kc-hist}

\subsubsection{Penyiapan Atribut dan Pelatihan Model MLPRegressor}
Kolom penanda non-prediktif (\texttt{id} dan \texttt{date}) beserta variabel target (\texttt{price}) dihilangkan dari matriks fitur masukan. Nilai fitur diproses melalui \texttt{SimpleImputer} berbasis rata-rata, lalu dibagi menjadi 70\% data latih dan 30\% data uji.

\screenshotimage[0.85\textwidth]{gambar/colab-p8-kc-cats.png}{Pemeriksaan Tipe Kolom pada df\_X KC House (Gambar 8.4.17 Modul)}{colab-kc-cats}

\begin{lstlisting}[style=pythonstyle, caption={Pelatihan Model MLPRegressor pada KC House Data}]
from sklearn.neural_network import MLPRegressor
from sklearn.model_selection import train_test_split
from sklearn.impute import SimpleImputer
import numpy as np

df_X = df.drop(['id', 'date', 'price'], axis=1)
df_y = df['price']

imputer = SimpleImputer(strategy='mean')
X = imputer.fit_transform(df_X)
y = df_y.astype(float).values

X_train, X_test, y_train, y_test = train_test_split(X, y, test_size=0.3, random_state=42)

reg = MLPRegressor(max_iter=500, random_state=42)
reg.fit(X_train, y_train)

print('coef of determination training :', reg.score(X_train, y_train))
print('coef of determination testing  :', reg.score(X_test, y_test))

y_pred = reg.predict(X_test)
print('prediction :', y_pred[:10])
print('real value :', y_test[:10])
\end{lstlisting}

\screenshotimage[0.88\textwidth]{gambar/colab-p8-kc-perf.png}{Skor Koefisien Determinasi dan Prediksi 10 Nilai Pertama (Gambar 8.4.18 Modul)}{colab-kc-perf}

\subsubsection{Pengukuran Metrik Performa Regresi ($R^2$ dan RMSE)}
Evaluasi regresi dihitung menggunakan metrik Root Mean Squared Error (RMSE) dan skor $R^2$.

\begin{lstlisting}[style=pythonstyle, caption={Pengukuran Metrik RMSE dan $R^2$}]
from sklearn.metrics import mean_squared_error
from sklearn.metrics import r2_score
import numpy as np

mse = mean_squared_error(y_test, y_pred)
rmse = np.sqrt(mse)
r2 = r2_score(y_test, y_pred)

print('rmse :', rmse)
print('r2   :', r2)
\end{lstlisting}

\screenshotimage[0.75\textwidth]{gambar/colab-p8-kc-eval.png}{Output Pengukuran Metrik RMSE dan R2 (Gambar 8.4.19 Modul)}{colab-kc-eval}

Hasil pengujian menghasilkan skor determinasi data latih sebesar 0,5636 dan data uji sebesar 0,5485, dengan nilai RMSE sebesar 255.319,80 USD.

\subsubsection{Visualisasi Komparasi Nilai Asli vs Prediksi}
Sebanyak 30 sampel data uji pertama dipetakan ke dalam diagram batang komparatif untuk melihat kedekatan antara estimasi model dan nilai riil properti.

\screenshotimage[0.92\textwidth]{gambar/colab-p8-kc-bar-pred.png}{Diagram Batang Perbandingan Nilai Riil vs Prediksi 30 Sampel (Gambar 8.4.20 Modul)}{colab-kc-bar-pred}

\clearpage

%======================================================================
% PEMBAHASAN TUGAS DAN ANALISIS
%======================================================================
\section{Pembahasan Tugas Praktikum}

\subsection{Tugas 1: Analisis Karakteristik dan Fitur Dataset}

\subsubsection{Tujuan Masing-Masing Dataset}
\begin{enumerate}
    \item \textbf{Bank Customer Churn Dataset (\texttt{BankChurners.csv})}:
    Dataset ini bertujuan memetakan perilaku pelanggan layanan kartu kredit perbankan guna mendeteksi dini nasabah yang memiliki probabilitas tinggi untuk menutup rekening atau berhenti berlangganan (\textit{churn}). Dalam operasional perbankan ritel, biaya akuisisi pelanggan baru jauh lebih mahal daripada biaya mempertahankan nasabah lama. Model prediktif yang dibangun dari dataset ini memberi peringatan otomatis bagi divisi manajemen relasi nasabah (\textit{Customer Relationship Management}) untuk menyalurkan program retensi khusus, penyesuaian limit kredit, maupun penawaran insentif bunga sebelum nasabah berpindah ke bank kompetitor.
    \item \textbf{King County House Sales Dataset (\texttt{kc\_house\_data.csv})}:
    Dataset ini mencatat rekam transaksi riil penjualan unit properti residensial di wilayah King County, Washington (mencakup metropolitan Seattle) sepanjang periode Mei 2014 hingga Mei 2015. Tujuannya adalah membangun model estimasi valuasi properti otomatis (\textit{Automated Valuation Model}/AVM). Model ini memetakan kontribusi matematis fitur struktural (luas bangunan, mutu konstruksi, jumlah kamar), karakteristik lingkungan (pemandangan panorama, tepian perairan), dan koordinat spasial geografis terhadap nilai transaksi pasar suatu rumah.
\end{enumerate}

\subsubsection{Kamus Data dan Karakteristik Fitur Dataset}
Karakteristik seluruh atribut penyusun kedua dataset dirangkum secara komprehensif pada Tabel \ref{fig:colab-t1-churn-features} dan Tabel \ref{fig:colab-t1-kc-features}.

\screenshotimage[0.95\textwidth]{gambar/colab-p8-t1-churn-features.png}{Kamus Data dan Peran Fitur Bank Customer Churn Dataset}{colab-t1-churn-features}

\screenshotimage[0.95\textwidth]{gambar/colab-p8-t1-kc-features.png}{Kamus Data dan Peran Fitur King County House Sales Dataset}{colab-t1-kc-features}

\clearpage

\subsection{Tugas 2: Eksperimen Model MLP pada Dataset Baru}

\subsubsection{Tugas 2A: Klasifikasi Penyakit Diabetes dengan MLPClassifier}
Dataset baru yang digunakan untuk permasalahan klasifikasi adalah \textit{Pima Indians Diabetes Dataset} (\texttt{diabetes.csv}). Dataset ini memuat 768 rekam medis pasien perempuan berketurunan Indian Pima dengan 8 indikator diagnostik klinis untuk memprediksi apakah pasien positif mengidap diabetes (\texttt{Outcome}=1) atau negatif (\texttt{Outcome}=0).

Model dibangun menggunakan arsitektur Multilayer Perceptron dengan konfigurasi dua lapisan tersembunyi \texttt{hidden\_layer\_sizes=(64, 32)}, fungsi aktivasi ReLU, optimizer Adam, dan diuji pada 231 sampel data uji independen.

\begin{lstlisting}[style=pythonstyle, caption={Pelatihan MLPClassifier pada Dataset Diabetes}]
# Split dan standarisasi
X_train_d, X_test_d, y_train_d, y_test_d = train_test_split(
    X_diab, y_diab, test_size=0.3, random_state=42, stratify=y_diab
)
scaler_diab = StandardScaler().fit(X_train_d)
X_train_d_s = scaler_diab.transform(X_train_d)
X_test_d_s = scaler_diab.transform(X_test_d)

# Model MLP Multi-layer (64, 32)
mlp_diab = MLPClassifier(hidden_layer_sizes=(64, 32), activation='relu', solver='adam', max_iter=500, random_state=42)
mlp_diab.fit(X_train_d_s, y_train_d)
\end{lstlisting}

\screenshotimage[0.60\textwidth]{gambar/colab-p8-t2-diab-cm.png}{Confusion Matrix Klasifikasi Diabetes Menggunakan MLPClassifier}{colab-t2-diab-cm}

\screenshotimage[0.72\textwidth]{gambar/colab-p8-t2-diab-loss.png}{Kurva Penurunan Nilai Loss Pelatihan MLP pada Dataset Diabetes}{colab-t2-diab-loss}

\screenshotimage[0.85\textwidth]{gambar/colab-p8-t2-diab-report.png}{Laporan Metrik Evaluasi Klasifikasi Diagnosa Diabetes}{colab-t2-diab-report}

Berdasarkan hasil pengujian pada Gambar \ref{fig:colab-t2-diab-report}:
\begin{enumerate}
    \item Model menghasilkan akurasi sebesar \textbf{68,83\%} dengan presisi kelas non-diabetes sebesar 0,78 dan recall sebesar 0,72.
    \item Pada kelas pasien diabetes (Outcome=1), model mencapai recall sebesar \textbf{0,62} (mendeteksi 50 pasien positif secara akurat dari total 81 kasus sebenarnya).
    \item Kurva loss pada Gambar \ref{fig:colab-t2-diab-loss} mengonfirmasi proses pembelajaran yang stabil, di mana nilai error menyusut secara konsisten dari kisaran 0,7 menuju di bawah 0,25.
\end{enumerate}

\subsubsection{Tugas 2B: Regresi Biaya Asuransi Kesehatan dengan MLPRegressor}
Dataset baru yang dipilih untuk regresi adalah \textit{Medical Cost Personal Dataset} (\texttt{insurance.csv}). Dataset ini mencakup 1.338 data individu dengan fitur usia, jenis kelamin, BMI, jumlah tanggungan anak, status perokok, wilayah tempat tinggal, serta besaran premi klaim asuransi kesehatan tahunan (\texttt{charges}).

Fitur kategorikal (\texttt{sex}, \texttt{smoker}, \texttt{region}) diproses menggunakan \textit{one-hot encoding}, dilanjutkan dengan standarisasi fitur. Arsitektur jaringan dikonfigurasi dengan dua lapisan tersembunyi yang lebih lebar, yaitu \texttt{(128, 64)}, untuk menangkap interaksi biaya kesehatan non-linear yang tajam akibat pengaruh konsumsi rokok dan obesitas.

\begin{lstlisting}[style=pythonstyle, caption={Pelatihan MLPRegressor pada Medical Cost Dataset}]
# One-hot encoding dan standard scaling
df_ins_encoded = pd.get_dummies(df_ins, columns=['sex', 'smoker', 'region'], drop_first=True)
X_ins = df_ins_encoded.drop('charges', axis=1).values
y_ins = df_ins_encoded['charges'].values

X_tr_i, X_te_i, y_tr_i, y_te_i = train_test_split(X_ins, y_ins, test_size=0.3, random_state=42)
scaler_i = StandardScaler().fit(X_tr_i)
X_tr_i_s = scaler_i.transform(X_tr_i)
X_te_i_s = scaler_i.transform(X_te_i)

# MLP Regressor multi-layer
mlp_ins = MLPRegressor(hidden_layer_sizes=(128, 64), activation='relu', solver='adam', max_iter=1000, random_state=42)
mlp_ins.fit(X_tr_i_s, y_tr_i)
\end{lstlisting}

\screenshotimage[0.72\textwidth]{gambar/colab-p8-t2-ins-scatter.png}{Scatter Plot Prediksi Model MLP vs Nilai Aktual Tagihan Asuransi}{colab-t2-ins-scatter}

\screenshotimage[0.92\textwidth]{gambar/colab-p8-t2-ins-bar.png}{Perbandingan Tagihan Asuransi Aktual vs Prediksi (30 Sampel Uji)}{colab-t2-ins-bar}

\screenshotimage[0.72\textwidth]{gambar/colab-p8-t2-ins-loss.png}{Kurva Konvergensi Loss Pelatihan MLPRegressor pada Biaya Asuransi}{colab-t2-ins-loss}

\screenshotimage[0.90\textwidth]{gambar/colab-p8-t2-ins-eval.png}{Rangkuman Kinerja dan Metrik Evaluasi MLPRegressor Tagihan Asuransi}{colab-t2-ins-eval}

Hasil eksperimen pada Gambar \ref{fig:colab-t2-ins-eval} memperlihatkan performa pemodelan regresi yang sangat prima:
\begin{enumerate}
    \item Model \texttt{MLPRegressor} mencatatkan koefisien determinasi $R^2$ sebesar \textbf{0,8437}. Artinya, arsitektur MLP mampu menjelaskan 84,37\% variansi besaran tagihan premi asuransi kesehatan individu.
    \item Nilai Root Mean Squared Error (RMSE) yang dicapai berada pada angka \textbf{4.787,68 USD} dengan Mean Absolute Error (MAE) sebesar \textbf{2.891,45 USD}.
    \item Scatter plot pada Gambar \ref{fig:colab-t2-ins-scatter} menunjukkan titik-titik data terdistribusi rapat di sepanjang garis diagonal ideal, membuktikan kemampuan jaringan syaraf tiruan dalam membedakan lonjakan biaya ekstrem antara nasabah bukan perokok dan nasabah perokok aktif secara presisi.
\end{enumerate}

\clearpage

%======================================================================
% BAB IV: KESIMPULAN
%======================================================================
\chapter[Kesimpulan]{Kesimpulan}

\begin{enumerate}
    \item Algoritma \textit{Multilayer Perceptron} (MLP) terbukti andal dalam menyelesaikan permasalahan pemodelan data non-linear, baik untuk tugas klasifikasi diskret menggunakan \texttt{MLPClassifier} maupun estimasi nilai numerik kontinu menggunakan \texttt{MLPRegressor}.
    \item Tahapan pra-pemrosesan data berupa \textit{categorical encoding} dan standarisasi skala fitur (\texttt{StandardScaler}) merupakan prasyarat mutlak dalam pelatihan neural network. Standarisasi mencegah terjadinya dominasi gradien oleh fitur berorde besar serta mempercepat konvergensi algoritma optimasi backpropagation.
    \item Pada klasifikasi nasabah perbankan (\texttt{BankChurners.csv}), model MLP berhasil mencapai akurasi sebesar 93,81\%, presisi 96,05\%, dan recall 96,58\%, serta mampu melakukan inferensi data nasabah baru secara akurat dengan tingkat keyakinan 96,52\%.
    \item Pada regresi harga properti King County (\texttt{kc\_house\_data.csv}), model MLPRegressor menghasilkan koefisien determinasi sebesar 0,5485 dengan RMSE 255.319,80 USD.
    \item Pada pengujian dataset baru (Tugas 2), arsitektur multi-layer \texttt{MLPClassifier} pada dataset diabetes menghasilkan akurasi 68,83\% dengan kemampuan mendeteksi 62\% kasus positif diabetes. Sementara itu, \texttt{MLPRegressor} pada dataset biaya asuransi kesehatan mencapai performa unggul dengan nilai $R^2$ sebesar 0,8437 dan RMSE 4.787,68 USD, membuktikan kapasitas representasi neural network dalam memetakan fungsi resiko finansial yang kompleks.
\end{enumerate}

\clearpage

%======================================================================
% DAFTAR PUSTAKA
%======================================================================
\addcontentsline{toc}{chapter}{Daftar Pustaka}
\begin{thebibliography}{99}

\bibitem{ref_goodfellow}
I.~Goodfellow, Y.~Bengio, and A.~Courville, \textit{Deep Learning}. Cambridge, MA, USA: MIT Press, 2016.

\bibitem{ref_haykin}
S.~Haykin, \textit{Neural Networks and Learning Machines}, 3rd~ed. Upper Saddle River, NJ, USA: Pearson Education, 2009.

\bibitem{ref_hukubun}
A.~J.~Hukubun, ``Neural Network,'' \textit{ResearchGate Technical Report}, 2022. Tersedia: \url{https://www.researchgate.net/publication/365581579_Neural_Network}.

\bibitem{ref_mardiani}
E.~Mardiani, N.~Rahmansyah, S.~Ningsih, D.~A.~Lantana, A.~S.~Wirawan, S.~A.~Wijaya, and D.~N.~Putri, ``Komparasi Metode KNN, Naive Bayes, Decision Tree, Ensemble, Linear Regression Terhadap Analisis Performa Pelajar SMA,'' \textit{INNOVATIVE: Journal Of Social Science Research}, vol.~3, no.~3, pp. 13880--13892, 2023.

\bibitem{ref_nielsen}
M.~A.~Nielsen, \textit{Neural Networks and Deep Learning: A Textbook}. Determination Press, 2015.

\bibitem{ref_pardede}
D.~Pardede, B.~H.~Hayadi, and Iskandar, ``Kajian Literatur Multi Layer Perceptron: Seberapa Baik Performa Algoritma Ini,'' \textit{Journal of ICT Application and System}, vol.~1, no.~1, pp. 23--34, 2022.

\bibitem{ref_roihan}
A.~Roihan, P.~A.~Sunarya, and A.~S.~Rafika, ``Pemanfaatan Machine Learning dalam Berbagai Bidang: Review Paper,'' \textit{IJCIT (Indonesian Journal on Computer and Information Technology)}, vol.~5, no.~1, pp. 75--82, 2020.

\end{thebibliography}

\end{document}
